% Section 1: Background and scope

\subsection{The Erd\H{o}s unit-distance problem}

For a finite set $P \subset \R^2$, let
$\nu(P) = \#\{\{x, y\} \subset P : |x - y| = 1\}$
and $\nu(n) = \max_{|P| = n} \nu(P)$. The planar unit-distance problem
\cite{Erdos1946} concerns the asymptotic behaviour of $\nu(n)$. Erd\H{o}s
proved $\nu(n) \geq n^{1 + \Omega(1/\log \log n)}$ via a square-grid
construction and conjectured the matching upper bound
$\nu(n) \leq n^{1 + C / \log \log n}$ for some absolute constant $C$.

\subsection{The 2026 disproof}

The conjecture was disproved in May 2026 by an internal reasoning model
at OpenAI \cite{OpenAI2026Proof}, with a simplified construction and
explicit numerical lower bound published the following day by Alon,
Bloom, Gowers, Litt, Sawin, Shankar, Tsimerman, Wang, and Matchett Wood
\cite{Sawin2026Remarks}. The combined result is:

\medskip
\noindent\textbf{Theorem 1.1 of \cite{Sawin2026Remarks}.}
\textit{There exists an absolute constant $\delta > 0$ and infinitely many
positive integers $n$ such that $\nu(n) \geq n^{1 + \delta}$.}

\medskip
The construction passes through a CM field $K_j = L_j(i)$, where the
$L_j$ form an infinite tower of totally real fields of bounded root
discriminant, supplied by the Golod--Shafarevich theorem
\cite{GolodShafarevich1964}. Section~2 of \cite{Sawin2026Remarks}
gives an explicit choice with
$T = \{3, 5, 7, 11, 13, 17\}$, $S = \{101, \infty\}$, and
$L_T = \Q(\sqrt{5}, \sqrt{13}, \sqrt{17}, \sqrt{21}, \sqrt{33})$,
and computes (eq.~(2.2)) a rigorous numerical lower bound
$\delta \geq \approx 6.24 \times 10^{-38}$.

\subsection{What this artifact verifies}
\label{sec:scope}

This paper accompanies a Python package, \texttt{erdos-ant-verification},
whose purpose is narrow and explicitly bounded:

\begin{itemize}[leftmargin=*]
  \item we \emph{do not} re-prove Theorem~1.1;
  \item we \emph{do not} construct the infinite Golod--Shafarevich tower;
  \item we \emph{do} re-implement the finite, explicitly checkable parts
        of the construction in pure Python, with exact small-integer
        arithmetic for the algebraic-number-theory layer and
        \texttt{mpmath} at 200-bit precision for the numerical evaluation
        of eq.~(2.2);
  \item we \emph{do} pin the published numerical value
        $\approx 6.24 \times 10^{-38}$ as a regression test that fails if
        our code drifts from the published source.
\end{itemize}

The artifact is therefore neither a proof checker (it does not formalise
the proof in a trusted kernel) nor an independent mathematical
contribution (it produces no new theorem). It is a reproducibility
artifact in the sense of, e.g., the ACM Artifact Evaluation guidelines:
a citable, version-pinned, CI-verified executable that allows any
third party to re-derive the published numerical value on their own
hardware without relying on our reading of the source PDFs.

The remainder of the paper describes the three implementation phases
(\S\S\ref{sec:phase1}--\ref{sec:phase3}), the central numerical
reproduction (\S\ref{sec:eq22}), the explicit limits
(\S\ref{sec:limits}), and the reproducibility surface
(\S\ref{sec:repro}).
